\hypertarget{conclusion-and-future-work}{%
\chapter{Conclusion and Future Work}\label{conclusion-and-future-work}}

\begin{draftnote}

\textbf{Draft status, 04 November 2026.} \S{}8.3 and \S{}8.6 depend on how the
project's claimed contribution is finally worded, which is pending a
conversation with the project guide. See the flags in Chapter 2 \S{}2.7 and
Chapter 7 \S{}7.9.

\end{draftnote}

\begin{center}\rule{0.5\linewidth}{0.5pt}\end{center}

\hypertarget{what-was-built}{%
\section{What was built}\label{what-was-built}}

A working three-tier system that accepts programming assignment
submissions, ranks them for human review, and attaches to each ranking
the evidence that produced it.

\begin{itemize}
\tightlist
\item
  \textbf{35 HTTP endpoints} --- 32 across four API routers, 3 on the
  detector.
\item
  \textbf{Eight document collections}, including an audit log held
  separately from the records it describes.
\item
  \textbf{Two detection layers in production}: structural similarity
  over normalised abstract syntax trees, and behavioural fingerprinting
  against a per-student baseline.
\item
  \textbf{A renormalised composition} of those layers into one Review
  Priority Score that stores the weights and threshold that produced it.
\item
  \textbf{A web interface} of ten pages, including a reviewer queue, an
  evidence view that guards its own arithmetic, and a student view that
  deliberately shows no score.
\item
  \textbf{278 automated tests}, all of them in the continuous
  integration gate since 01 November 2026.
\item
  \textbf{A committed, re-runnable measurement programme} of seven
  scripts.
\end{itemize}

\hypertarget{what-the-project-established}{%
\section{What the project
established}\label{what-the-project-established}}

\textbf{The architecture works and is defensible.} Credentials exist in
exactly one tier; an absent layer contributes nothing rather than zero;
a result remains interpretable after the configuration that produced it
has changed; a re-run supersedes without destroying. None of these is
remarkable individually, and together they are what makes the output
auditable rather than merely computed.

\textbf{Layer 1 detects copying, measured.} In the positive control of
the ablation study, with a donor file inside the comparison pool, Layer
1's separation of spliced from clean files rises monotonically with how
much was copied --- +0.0150, +0.0335, +0.0706 --- significant at all
three fractions, with precision@10 reaching 1.000 where half the file
was taken.

\textbf{The style features carry some author signal, within a project.}
Across eleven within-project tests, roughly 1.4\texttimes{} to 2.5\texttimes{} over chance
where they work at all --- and they fail to beat chance in more than
half the projects tested. That is the honest figure and it survives
checking.

\textbf{The two layers are independent.} Pearson -0.042 across 428
scored files, so the premise on which the RPS weights them separately
holds.

\textbf{The measurement programme is itself a result.} Every figure in
Chapter 7 is produced by a committed script that can be re-run, several
of them reporting ranges across 25 draws rather than point estimates ---
a practice adopted because a single sample proved unreliable, not
because it was planned.

\hypertarget{what-the-project-did-not-establish}{%
\section{What the project did not
establish}\label{what-the-project-did-not-establish}}

\textbf{Layer 2 does not detect partial outsourcing, and that is the
project's headline negative result.} On the 214-pair style-shift corpus
--- 94 independent source files --- the behavioural layer did not
separate spliced submissions from clean ones at any splice fraction, and
combining it with Layer 1 degraded ranking precision.

The mechanism is understood and is structural rather than a matter of
tuning, and every link in the explanation is measured rather than
argued:

\begin{enumerate}
\def\labelenumi{\arabic{enumi}.}
\tightlist
\item
  The donors genuinely differ from the targets, by about 0.9 cohort
  standard deviations.
\item
  The \textbf{direction} of that difference is arbitrary: the donor sits
  above the target in 0.462 to 0.500 of pairs across the four features,
  none of them distinguishable from a coin flip.
\item
  The splice realises the predicted displacement almost exactly, 0.89 to
  1.00.
\end{enumerate}

A score that is a \textbf{magnitude} of deviation from a baseline
therefore cannot see it: a displacement of predicted size and arbitrary
sign, applied to a point already off centre, moves it closer about as
often as further away. Separately, \texttt{max\_block\_depth} is nearly
inert for this task because a maximum is not a mixable statistic.

\textbf{Layer 1's precision and recall on genuine plagiarism are
unknown}, because no labelled corpus of real student plagiarism exists.
What was measured is the distribution over honest work by different
authors, which is a different claim.

\hypertarget{future-work}{%
\section{Future work}\label{future-work}}

Ordered by what unblocks the most, with preconditions stated rather than
implied.

\textbf{1. Redesign Layer 2 around per-unit features and a signed
comparison.} This follows directly from \S{}8.3 and is the single most
valuable item. Two changes are needed together: features computed per
function or per window, so a spliced region can be isolated instead of
averaged into its host; and a signed or population-referenced
comparison, so that moving \emph{toward} another author's typical values
is itself the signal rather than being invisible to a distance.
\emph{Precondition: none. This is buildable now and the corpus to test
it already exists.}

\textbf{2. Implement Layer 3 for machine-generated code.}
\emph{Precondition: a labelled corpus of student submissions from the
course the system will police.} Without it the layer can be built but
never characterised, and an uncharacterisable layer in an evidence
system is worse than an absent one (Chapter 5 \S{}5.4). The precondition is
the work, not the model.

\textbf{3. Reconsider the scored feature set.} The current four were
selected from a 60-file sample and that selection is retracted: a
204-file measurement found that nine of ten features clear F = 1 in
Python, that the feature which drove the original choice ranks ninth,
and that using every applicable feature beats the fixed four in both
languages. A fixed four also cannot be right for both, since two of them
behave oppositely across the languages. \emph{Precondition: none, but it
should follow item 1 rather than precede it, since the aggregation level
is the larger problem.}

\textbf{4. Build a labelled corpus of genuine plagiarism.} This is what
would turn Layer 1's distributional measurement into a precision and
recall figure. \emph{Precondition: institutional access to adjudicated
cases, with consent and anonymisation --- which is an ethics process,
not an engineering task.}

\textbf{5. The deferred interface work.} WebSocket push for the review
queue, the side-by-side AST diff, and multi-file or archive submissions.

\textbf{6. Extend the ablation to Java.} The present study is Python
only, and the two most affected features behave oppositely across the
two languages, so nothing should be assumed to transfer.

\hypertarget{on-method}{%
\section{On method}\label{on-method}}

Three practices account for most of what this project got right, and all
three are transferable.

\textbf{Measure, then retract when the measurement disagrees.} Several
published claims in this project were withdrawn because a larger or
better-controlled measurement contradicted them: a similarity floor that
did not hold, a feature ranking that did not replicate, an AUC pair
traceable to no run. Each retraction is recorded with its reasoning
rather than quietly corrected.

\textbf{Read code as well as run it.} Four defects in the ablation study
were found by reading, and not one would have been caught by running it,
because each produced output that looked like a result --- a control
that compared the wrong file, a chance baseline that inflated every
figure, a measurement of a strategy the system does not deploy, and a
significance verdict printed without its direction.

\textbf{Render data and look at it.} Nine defects were found by
displaying real data on a screen and reading it, and two more by
printing it to paper. A test asserts a value; every one of these was
something nobody had thought to assert.

\hypertarget{closing}{%
\section{Closing}\label{closing}}

The system set out to rank submissions for three kinds of misconduct:
collusion, outsourcing and machine generation. It ranks reliably for the
first, has a documented and justified absence for the third, and for the
second it produced a measured negative result with a diagnosed mechanism
and a named fix.

\begin{draftnote}

\textbf{To be finalised with the project guide.} How this chapter states
the project's contribution depends on how the claimed novelty is
reworded in light of \S{}8.3. A negative result with a measured mechanism
and a specific design remedy is a defensible final-year outcome; a claim
contradicted by the project's own committed measurement script is not.
Chapters 2 and 7 carry the same flag, and all three should be settled in
one conversation.

\end{draftnote}
